\section{Further discussion}
\label{sec:discussion}

This section expands upon the main findings of the systematic literature review by offering critical reflections and contextual analysis. The aim is to highlight cross-cutting insights and to examine how methodological and contextual differences influence the interpretation and application of churn prediction research.

\subsection{Churn definitions}
\noindent An analysis of the selected studies reveals a notable diversity in the definitions of customer churn. This variation suggests that churn definitions are often tailored to the specific characteristics of the retail context, including the type of retailer (e.g., e-commerce vs. brick-and-mortar) and the business sector in which they operate. 
%It may, in fact, be beneficial for businesses of the same retail type but operating in different market segments to adopt distinct churn definitions, due to differences in customer behavior and distribution.
Even among businesses of the same retail type, adopting distinct churn definitions may be appropriate, especially when customer behavior and distribution patterns differ across market segments.
For instance, consider two grocery retailers: one located in a highly touristic area and the other in a more residential setting with stable purchasing patterns. In the former case, customer behavior is likely to be more variable and influenced by seasonality or transient visits, whereas in the latter, patterns may be more predictable and consistent. A fixed definition of churn (e.g., no purchases within a specific time window) may not perform equally well across both scenarios. 
%In such situations, it is more effective to adopt a dynamic, customer-specific definition of churn that takes into account the individual’s purchase history, resulting in more accurate labeling and improved model performance. 
Dynamic, customer-specific definitions that account for individual purchase history are likely to result in more accurate labeling and improved model performance in such contexts.
Similarly, differences across retail sectors may also necessitate tailored churn definitions, as a single approach may not be equally suitable for all types of businesses.

\subsection{AI techniques employed for churn prediction}

\noindent While a variety of AI techniques have been applied to churn prediction, many studies continue to rely on conventional machine learning methods, such as decision trees, logistic regression, and random forests. 
%It is possible that the performance landscape would shift if more modern architectures, such as transformer-based models or deeper neural networks, were more widely adopted.
The broader adoption of more modern architectures, such as transformer-based models or deep neural networks, could potentially alter the performance landscape. Given their ability to capture complex and nonlinear relationships between features, the relative importance of different variables in predicting churn could also change.

However, drawing definitive conclusions about the best-performing technique remains challenging due to several factors. First, the studies analyzed utilize different datasets, each with varying structures, feature sets, and levels of imbalance. Second, the definition of churn itself is not standardized and varies across studies: some define churn as complete inactivity, while others consider a decline in engagement within a specific timeframe. These inconsistencies hinder direct performance comparisons.

%Moreover, for the sake of analysis, techniques have been grouped into broad methodological clusters (e.g., neural networks, tree-based models). However, even within a given cluster, there may be substantial variation in model complexity, hyperparameter settings, and architectural design. 
Although we grouped techniques into broad methodological clusters (e.g., neural networks, tree-based models), substantial variation exists within each cluster—ranging from model complexity and hyperparameter settings to architectural design. 
For instance, studies that apply artificial neural networks may differ widely in terms of the number of layers, node configurations, activation functions, and training regimes employed. As such, the results should be interpreted with caution, recognizing that similar-sounding methods may vary significantly in implementation and performance.


\subsection{E-commerce vs Brick and Mortar Retailers}\nopagebreak
\noindent Research in customer churn prediction is more prevalent in the e-commerce domain, largely due to the availability of accessible and well-structured online datasets, including publicly available sources such as those hosted on Kaggle. E-commerce platforms benefit from extensive digital footprints left by users, which enable advanced predictive modeling using sophisticated machine learning techniques, including deep neural networks and transformer-based architectures. The digital environment inherently supports continuous and fine-grained tracking of customer interactions, leading to more accurate and timely churn predictions.

In contrast, churn prediction in brick-and-mortar retail---particularly in sectors such as fast-moving consumer goods (FMCG) and grocery sectors---relies primarily on transactional data. The nature and granularity of available data in physical retail environments are more limited. While loyalty programs allow retailers to link purchases to individual customers, a portion of in-store transactions remains anonymous and cannot be attributed to specific customer profiles. As a result, the effectiveness of churn prediction in these settings depends heavily on the existence and adoption rate of robust loyalty systems.

In e-commerce, the requirement for user account creation allows for consistent tracking of individual customer behavior, extending beyond transactions to include web navigational behavior, such as pages visited, time spent on site, product views, browsing frequency, and revisit intervals. These non-transactional behavioral indicators offer early signals of churn intent, enabling proactive retention strategies.
%businesses to act proactively and deploy retentive strategies before the customer disengages completely.

By contrast, such behavioral data is largely inaccessible in physical retail. This reinforces the comparative richness of e-commerce data, which enables the integration of both transactional and behavioral signals in churn modeling. 
It is important to note, however, that retailers with both physical and online channels, often referred to as omnichannel retailers, can improve the accuracy of churn prediction by integrating data from both sources into their data warehouse. 
Interestingly, some of the studies that were focused on e-commerce did not utilize web navigational data, despite operating in a digital context where such features are readily available and potentially valuable.


%Among the 41 studies reviewed, 25 were focused on e-commerce, while the rest addressed churn prediction in brick-and-mortar settings. Interestingly, some e-commerce studies (e.g., S5, S8, and S13) did not utilize web navigational data, despite operating in a digital context where such features are readily available and potentially valuable.

Ultimately, churn prediction serves as a strategic tool to optimize marketing efforts, reduce customer attrition, and support revenue retention. Businesses---especially in data-rich environments like e-commerce---can closely align churn prediction with marketing strategies to enable timely and personalized retention interventions, maximizing the impact of promotional resources.

\subsection{Final considerations}\nopagebreak
\noindent
\textcolor{blue}{Synthesizing the results across our RQs provides a clearer picture of how AI is currently applied to churn prediction in non-contractual retail and why evidence accumulation remains difficult. First, most studies implement a largely standard predictive pipeline: churn is operationalized through study-specific labeling rules (Section~\ref{sec:RQ1}), input signals are strengthened via engineered customer summaries (feature engineering is reported in 34 of 41 studies; Section~\ref{sec:RQ3}) and performance is then benchmarked across a set of supervised learners where tree-based models and neural networks frequently emerge as strong candidates but no single approach dominates across contexts (Section~\ref{sec:RQ4}). Second, the field’s implicit ``framework'' is therefore predominantly methodological and data-driven rather than protocol-driven: model choices are strongly conditioned by dataset availability and by the adopted churn definition, which limits cross-study comparability. This is reinforced by the scarcity of reusable benchmarks (only 11 of 41 studies provide publicly available datasets; Section~\ref{sec:RQ6}) and by the recurring single-source, context-specific settings emphasized among reported limitations (15 of 41 studies; Section~\ref{sec:RQ7}), which together constrain domain generalizability across retail sectors and channels. Finally, the gap between technical results and actionable retention decisions is visible in evaluation practice: the literature relies predominantly on standard classification metrics (Section~\ref{sec:RQ5}), while future directions explicitly call for explainability, ROI-aware assessment, and real-world validation to connect predictive accuracy to intervention cost and business impact (Section~\ref{sec:RQ8}). Taken together, these insights suggest that strengthening the paper’s contribution is less about proposing yet another algorithmic variant and more about enabling transferable evidence through data/domain augmentation, clearer reporting of churn operationalization, and evaluation designs aligned with deployment and ROI constraints.}
